\documentclass[]{article}

\usepackage{amsmath}
\usepackage{multirow}
\usepackage{natbib}
\usepackage{graphicx} 


\begin{document}

% Auto-generated figures file
% \usepackage{graphicx}

\begin{figure}[h!]
    \centering
    \includegraphics[width=0.5\textwidth]{../input_files/plots/step_3_analysis_plot_1777372490.png}
    \caption{Properties of the GEMS (blue points) and Sparkler (orange points) globular cluster (GC) samples. Panels (a) and (b) show the cumulative formation history from the Trenti et al. (2015) atomic cooling model. Panel (c) displays the mass-metallicity relation, where GEMS GCs show a weak positive correlation. Panel (d) plots metallicity against formation redshift ($z\_{form}$), revealing that GEMS GCs formed earlier ($z\_{form} \approx 10-19$) yet are more metal-rich than the Sparkler GCs ($z\_{form} \approx 2-4$). Panel (e) shows the age-mass relation for the GEMS sample. Panel (f) presents the age distributions, demonstrating that GEMS GCs are observed at very young ages (< 310 Myr) while Sparkler GCs are significantly older (> 1.5 Gyr).}
    \label{fig:image0}
\end{figure}



\end{document}
                